\documentclass[11pt,a4paper]{article}
\usepackage[margin=25mm]{geometry}
\usepackage{amsmath,amssymb,amsthm,booktabs,hyperref}
\setlength{\emergencystretch}{2em}
\newtheorem{theorem}{Theorem}
\newtheorem{lemma}[theorem]{Lemma}
\newcommand{\gen}[1]{\langle#1\rangle}
\newcommand{\Sym}{\operatorname{Sym}}
\title{Symmetric wings at a perfect split}
\author{Human Compression project: a component for L019}
\date{}
\begin{document}\maketitle
A string C-marking is an ordered family of nonidentity involutions
whose nonadjacent entries commute and which satisfies the intersection
property for every pair of subsets of its entries. Throughout, groups
are actual permutation groups. A group supported on $A$ fixes the
complement of $A$ pointwise.

\begin{theorem}\label{wing}
Let $\rho_0,\ldots,\rho_{r-1}$ be a transitive string C-marking on
the disjoint union $\Omega=A\sqcup B$, of size $n$. Suppose that,
for some $i$, the marks before $i$ are supported on $A$, the marks
after $i$ are supported on $B$, and
\[
 \rho_i=(a\ b)\alpha\beta,\qquad a\in A,\quad b\in B,
\]
where $\alpha$ and $\beta$ are involutions or identities supported
on $A\setminus\{a\}$ and $B\setminus\{b\}$, respectively.
Suppose also that the omission of $\rho_i$ has exactly the two
orbits $A,B$. Set
\[
 L=\gen{\rho_0,\ldots,\rho_{i-1}},\qquad
 R=\gen{\rho_{i+1},\ldots,\rho_{r-1}}.
\]
If $2r>n+2$, then $L=\Sym(A)$ or $R=\Sym(B)$.
The empty group on a singleton is understood as its full symmetric
group. More precisely, each wing which is not full symmetric has
rank at most half its degree.
\end{theorem}

These are the actual support hypotheses of a perfect split in
Cameron--Fernandes--Leemans, \emph{The number of string C-groups of
high rank}, arXiv:2212.12723v2. The theorem applies, in particular,
at their threshold $r\ge(n+3)/2$. Its proof does not assert that an
arbitrary restriction or block quotient of a C-marking is again a
C-marking.

\paragraph{Acquired classification inputs.}
We use three complete earlier recoveries, with their original
attributions and finite premises retained in the dependency register.
The transitive rank theorem of Cameron--Fernandes--Leemans--Mixer
(L047) classifies every faithful transitive string C-marking of
degree $d$, rank greater than $d/2$, and group different from
$A_d,S_d$. Its exceptions are the odd wreath family in
Lemma~\ref{nonalt}, the complete degree-six/eight imprimitive
markings, and the primitive degree-six rank-four $S_5$ marking.
The complete alternating rank theorem of the same four authors
(L044) leaves only $A_5$ of rank three and $A_{11}$ of rank six
above half the natural degree. Finally, the acquired alternating
census from L071 supplies all four ordered rank-six $A_{11}$
markings, up to simultaneous point conjugacy. Its complete
enumeration, including its bound on the searched ranks and its
intersection checks, is retained; it is not inferred from a list
of examples. We only perform the additional finite tests stated
below on these already complete classifications.

\begin{lemma}\label{anchor}
Let a nontrivial wing in Theorem~\ref{wing} have degree $d$ and
write its marks as $t_0,\ldots,t_{k-1}$ in order starting nearest
the split. Its attachment point $a$ is moved by $t_0$ and fixed
by all $t_j$, $j>0$. If the adjacent split mark is written
$q=(a\ b)\alpha\beta$ as above, then
\[
 \alpha\in C_{\Sym(A)}(t_1,\ldots,t_{k-1})_a.
\]
\end{lemma}
\begin{proof}
The omission is the direct product of its two supported wings.
Its two stated orbits therefore make each wing transitive on its
own support. The actual wing marks inherit the intersection
property as a subfamily; their restriction to their support is
faithful, so this changes no group or intersection.
For $j>0$, the string relation gives $t_jq=qt_j$. Since $t_j$
fixes $b$ and $q(b)=a$, it fixes $a$. Transitivity and $|A|>1$
then force $t_0$ to move $a$. Both $(a\ b)$ and $\beta$ commute
with $t_j$, so $\alpha$ does too. Its definition already fixes $a$.
\end{proof}

Call a point moved by the first mark and fixed by every later mark
an \emph{endpoint anchor}. The following argument inventories all
nonalternating exceptions, including the actual action before
any change by an abstract automorphism.

\begin{lemma}\label{nonalt}
No nonalternating proper transitive exception above half the degree
in L047 has an endpoint anchor, in either orientation.
\end{lemma}
\begin{proof}
For the uniform family the degree is $2m$, with $m\ge5$ odd,
and the rank is $m+1$. The acquired equality proof gives, in the
actual two-point block action, the two distinguished base elements
$\lambda$ and $\lambda+e_j$. Here $\lambda$ flips every block,
and $\lambda+e_j$ flips every block except $j$; the remaining
marks act transitively on the blocks. This conclusion precedes
the cocycle automorphism used to put their lifts in a convenient
abstract normal form.

If a point were moved by just one distinguished mark, that mark
would have to be $\lambda$. It would also have to lie in block
$j$, since otherwise $\lambda+e_j$ moves it. But some remaining
mark moves block $j$, because their block action is transitive
and $m>1$. It moves both its points, a contradiction. This
argument permits all lifts of the block marks and all allowed
string orders. The case $m=3$ is included in the finite census.

There are fourteen retained ordered imprimitive exception tuples:
thirteen of degree six and one of degree eight. The primitive
exception contributes one further degree-six tuple. For each
literal tuple $(t_0,\ldots,t_{k-1})$ the certificate evaluates
\[
 \{a:t_0(a)\ne a,\ t_j(a)=a\text{ for every }j>0\}
\]
and the same expression for its reverse. All thirty sets are
empty. The input is extracted exactly from the complete L047
census, with its indices and hashes retained. Evaluating these
sets is an exhaustive point test, and proves the assertion for
every tuple in that classification.
\end{proof}

\begin{lemma}\label{alternating}
Neither $A_5$ of rank three nor $A_{11}$ of rank six can be a
wing in Theorem~\ref{wing}.
\end{lemma}
\begin{proof}
In $A_5$, every nonidentity involution has a unique fixed point.
For a proposed anchor $a$ of $(t_0,t_1,t_2)$, that point for
$t_2$ would be $a$. Since $t_0$ commutes with $t_2$, it preserves
its fixed-point set and fixes $a$, a contradiction.

For $A_{11}$ the four complete ordered census representatives
share the marks
\[
\begin{aligned}
 t_0&=(8\ 9)(10\ 11),& t_1&=(7\ 8)(9\ 10),\\
 t_3&=(3\ 4)(5\ 6),& t_5&=(1\ 2)(3\ 4),
\end{aligned}
\]
and have the following choices. ``First'' means the displayed
orientation, and ``last'' means its reverse.
\begin{center}
\begin{tabular}{ccccc}
\toprule
Row&$t_2$&$t_4$&First anchors&Last anchors\\
\midrule
1&$(6\ 7)(10\ 11)$&$(2\ 3)(4\ 5)$&$\varnothing$&$\{1\}$\\
2&$(6\ 7)(10\ 11)$&$(1\ 3)(2\ 5)$&$\varnothing$&$\varnothing$\\
3&$(6\ 7)(8\ 9)$&$(2\ 3)(4\ 5)$&$\{11\}$&$\{1\}$\\
4&$(6\ 7)(8\ 9)$&$(1\ 3)(2\ 5)$&$\{11\}$&$\varnothing$\\
\bottomrule
\end{tabular}
\end{center}
For each of the four nonempty anchor cases, orient the tuple
with the anchor first, and denote it again by $(t_0,\ldots,t_5)$.
The complete centralizer fixing its anchor is
\[
 C_{S_{11}}(t_1,\ldots,t_5)_a=1.                 \tag{1}
\]
This is checked by an exhaustive equivariant-map calculation:
the image of a point determines the images of its whole orbit
under the five specified permutations. Start by fixing $a$;
at each still unassigned point try every unused possible image,
propagate through all five marks, and reject a conflicting image
or repeated value. Every centralizing bijection occurs in exactly
one branch. The only surviving bijection in each case is the
identity. Thus (1) is a complete centralizer computation, rather
than a test of selected involutions. Lemma~\ref{anchor} forces
$\alpha=1$.

Set $u=(a\ 12)$ and consider the first five entries
$(u,t_0,t_1,t_2,t_3)$ of the pure attachment. Direct finite closure
under right multiplication gives, in every anchor case,
\[
\begin{array}{c|c}
\text{subgroup}&\text{order}\\ \hline
 U=\gen{t_0,t_1,t_2,t_3}&5040\\
 V=\gen{u,t_0,t_1,t_2}&5040\\
 U\cap V&360\\
 M=\gen{t_0,t_1,t_2}&60.
\end{array}                                                   \tag{2}
\]
The literal-permutation certificate enumerates these finite
closures and intersects the two sets. A separate GAP computation
on the exact retained census agrees. In particular
$[U\cap V:M]=6$.

In the actual attachment $q=u\beta$, the permutation $\beta$
has disjoint support from all five displayed pure marks. The
group on the corresponding actual entries is a subgroup of
$\gen{u,t_0,t_1,t_2,t_3}\times\gen\beta$. Projection onto the
first factor has central kernel of order at most two. The
acquired central-intersection lemma says that, if its marking
had the intersection property, the index in (2) would be a
power of two: the quotient of the downstairs end intersection
by the common-middle image embeds in a quotient of that central
kernel. The factor three in six is a contradiction. This also
covers $\beta=1$. As a submarking of a C-marking must satisfy
the intersection property, no such wing occurs.
\end{proof}

\begin{proof}[Proof of Theorem~\ref{wing}]
If either support is a singleton, its supported wing is its full
symmetric group. Otherwise both wings are faithful transitive
string C-groups with the endpoint anchors of Lemma~\ref{anchor}.
For a proper nonsymmetric wing, the complete L047 theorem and
Lemma~\ref{nonalt} exclude every nonalternating rank above half
the degree. The complete L044 theorem and
Lemma~\ref{alternating} exclude every alternating rank above half
the degree. Consequently, if both wings were proper, their ranks
would satisfy
\[
 r-1=i+(r-i-1)\le |A|/2+|B|/2=n/2.
\]
This contradicts $2r>n+2$ and proves the theorem.
\end{proof}

\paragraph{Attribution, computation and scope.}
The three classification theorems remain attributed to their
historical authors. Their full acquired proofs and computational
premises are dependencies, not new results of this note. The new
argument combines their exact exception lists with actual endpoint
anchors, the full anchored centralizers, and the central-intersection
obstruction. The central-intersection lemma is the separately reviewed
component of this L019 recovery. The finite tests here use literal
permutation arithmetic independently of GAP; GAP supplies a second
check and the earlier complete classifications retain their stated
software and library premises. No new census is claimed.

This proves the stated perfect-split wing theorem only. It neither
establishes the existence of a perfect split nor completes the
extension/contraction bijection or the eventual counting constants
of L019. It makes no assertion that every auxiliary proposition
in the source is valid, and it grants no whole-original recovery,
open-conjecture or Lean credit.
\end{document}
